\documentclass{article}
\usepackage{amsmath,amssymb,amsthm}
\usepackage{algorithm}
\usepackage{algpseudocode}
\usepackage{enumitem}
\usepackage{hyperref}
\usepackage{bm}
\usepackage{graphicx}
\usepackage{color}
\usepackage{mathtools}

\newtheorem{theorem}{Theorem}
\newtheorem{lemma}{Lemma}
\newtheorem{assumption}{Assumption}

\begin{document}

%%%%%%%%%%%%%%%%%%%%%%%%%%%%%%%%%%%%%%%%%%%%%%%%%%%%%%%%%%%%%%%%%%%%%%%%%%%%
\section*{Clipped Stochastic Gradient Descent (Clipped SGD)}
%%%%%%%%%%%%%%%%%%%%%%%%%%%%%%%%%%%%%%%%%%%%%%%%%%%%%%%%%%%%%%%%%%%%%%%%%%%%

\section*{Mathematical Formulation}
Let $f:\mathbb{R}^d\rightarrow\mathbb{R}$ be the objective function and let
$\{z_t\}_{t\ge 1}$ denote an i.i.d.\ data stream.
For each iteration $t$ we compute a (possibly stochastic) gradient
\[
g_t^{\text{raw}} \;:=\; \nabla f(x_t;z_t)\in\mathbb{R}^d .
\]
Given a clipping threshold $C>0$, the \emph{clipped gradient} is defined as
\begin{equation}
\label{eq:clipped}
g_t \;:=\;
\frac{g_t^{\text{raw}}}{\displaystyle \max\!\Bigl\{1,\,
\frac{\|g_t^{\text{raw}}\|}{C}\Bigr\}}
\;=\;
\begin{cases}
g_t^{\text{raw}}, & \text{if }\|g_t^{\text{raw}}\|\le C,\\[4pt]
C\,\dfrac{g_t^{\text{raw}}}{\|g_t^{\text{raw}}\|}, & \text{if }\|g_t^{\text{raw}}\|> C .
\end{cases}
\end{equation}
With a stepsize (learning rate) schedule $\{\eta_t\}_{t\ge 1}$ the update rule is
\begin{equation}
\label{eq:update}
x_{t+1}=x_t-\eta_t\, g_t .
\end{equation}
The algorithm stops after $T$ iterations or when a prescribed tolerance
on the (clipped) gradient norm is reached.

\section*{Pseudocode}
\begin{algorithm}[H]
\caption{Clipped SGD}
\begin{algorithmic}[1]
\State \textbf{Input:} stepsizes $\{\eta_t\}_{t=1}^T$, clipping threshold $C>0$, initial point $x_1\in\mathbb{R}^d$.
\For{$t=1,\dots,T$}
    \State Sample data $z_t$ (or draw a minibatch).
    \State Compute raw gradient $g_t^{\text{raw}} \gets \nabla f(x_t;z_t)$.
    \State Compute clipping factor 
    \[
        \alpha_t \gets \max\!\Bigl\{1,\frac{\|g_t^{\text{raw}}\|}{C}\Bigr\}.
    \]
    \State Form clipped gradient $g_t \gets g_t^{\text{raw}}/\alpha_t$.
    \State Update $x_{t+1} \gets x_t - \eta_t\, g_t$.
\EndFor
\State \textbf{Output:} $x_{T+1}$ (or the average iterate $\bar x_T=\frac1T\sum_{t=1}^T x_t$).
\end{algorithmic}
\end{algorithm}

\section*{Dry Run Example}
Consider the simple one–dimensional quadratic
\[
f(w)=(w-3)^2 ,
\qquad \nabla f(w)=2(w-3).
\]
Set $C=1$, stepsize $\eta=0.1$, and initialise $w_0=0$.

\begin{center}
\begin{tabular}{c|c|c|c|c}
$t$ & $w_t$ & $\nabla f(w_t)$ & $\| \nabla f(w_t)\|$ & $g_t$ (clipped) \\ \hline
0 & $0$ & $-6$ & $6> C$ & $g_0 = C\cdot\frac{-6}{6}= -1$ \\[4pt]
1 & $w_1 = 0-0.1(-1)=0.1$ & $2(0.1-3)=-5.8$ & $5.8> C$ & $g_1 = -1$ \\[4pt]
2 & $w_2 = 0.1-0.1(-1)=0.2$ & $2(0.2-3)=-5.6$ & $5.6> C$ & $g_2 = -1$ \\[4pt]
3 & $w_3 = 0.2-0.1(-1)=0.3$ & $2(0.3-3)=-5.4$ & $5.4> C$ & $g_3 = -1$ 
\end{tabular}
\end{center}
The objective values are
\[
f(w_0)=9,\quad f(w_1)= (0.1-3)^2 = 8.41,\quad
f(w_2)=8.04,\quad f(w_3)=7.69 .
\]
Even though the true gradient magnitude is far larger than $C$, clipping to $C$
produces a stable, bounded update direction.

%%%%%%%%%%%%%%%%%%%%%%%%%%%%%%%%%%%%%%%%%%%%%%%%%%%%%%%%%%%%%%%%%%%%%%%%%%%%
\section*{Assumptions}
%%%%%%%%%%%%%%%%%%%%%%%%%%%%%%%%%%%%%%%%%%%%%%%%%%%%%%%%%%%%%%%%%%%%%%%%%%%%
\begin{assumption}[Smoothness]\label{ass:smooth}
$f$ is $L$‑smooth, i.e. for all $x,y\in\mathbb{R}^d$,
\[
\|\nabla f(x)-\nabla f(y)\|\le L\|x-y\|\quad\Longleftrightarrow\quad
f(y)\le f(x)+\langle\nabla f(x),y-x\rangle+\frac{L}{2}\|y-x\|^2 .
\]
\end{assumption}
\begin{assumption}[Lower Boundedness]\label{ass:lb}
There exists $f_{\inf}>-\infty$ such that $f(x)\ge f_{\inf}$ for all $x$.
\end{assumption}
\begin{assumption}[Unbiased Stochastic Gradient]\label{ass:unbiased}
For every $t$, $\mathbb{E}[g_t^{\text{raw}}\mid x_t]=\nabla f(x_t)$.
\end{assumption}
\begin{assumption}[Bounded Variance]\label{ass:variance}
There exists $\sigma^2\ge0$ with
\[
\mathbb{E}\bigl[\|g_t^{\text{raw}}-\nabla f(x_t)\|^2\mid x_t\bigr]\le\sigma^2,
\qquad\forall t .
\]
\end{assumption}
\begin{assumption}[Stepsize]\label{ass:stepsize}
The stepsizes satisfy $\eta_t=\eta$ (constant) with
$0<\eta\le \frac{1}{L}$.
\end{assumption}
All constants $L,C,\sigma,\eta$ are known a‑priori.

%%%%%%%%%%%%%%%%%%%%%%%%%%%%%%%%%%%%%%%%%%%%%%%%%%%%%%%%%%%%%%%%%%%%%%%%%%%%
\section*{Main Convergence Theorem}
%%%%%%%%%%%%%%%%%%%%%%%%%%%%%%%%%%%%%%%%%%%%%%%%%%%%%%%%%%%%%%%%%%%%%%%%%%%%
\begin{theorem}[Convergence of Clipped SGD]\label{thm:main}
Under Assumptions \ref{ass:smooth}–\ref{ass:stepsize}, after $T$ iterations
\[
\frac{1}{T}\sum_{t=1}^{T}
\mathbb{E}\!\Bigl[\min\!\bigl\{C,\;\|\nabla f(x_t)\|\bigr\}^2\Bigr]
\;\le\;
\frac{2\bigl(f(x_1)-f_{\inf}\bigr)}{\eta T}
\;+\;
2L\eta\,\sigma^2 .
\]
Consequently, choosing $\eta=\Theta\bigl(\tfrac{1}{\sqrt{T}}\bigr)$ yields
\[
\frac{1}{T}\sum_{t=1}^{T}
\mathbb{E}\!\Bigl[\min\!\bigl\{C,\;\|\nabla f(x_t)\|\bigr\}^2\Bigr]
= O\!\bigl(T^{-1/2}\bigr).
\]
\end{theorem}

\section*{Theoretical Properties}
\begin{itemize}[leftmargin=*]
\item \textbf{Stability.} Because $ \|g_t\|\le C$ by construction, the update
$magnitude$ $\eta C$ is uniformly bounded, preventing exploding steps even
when raw gradients have arbitrarily large variance.
\item \textbf{Bias induced by clipping.}
When $\|\nabla f(x_t)\|>C$ the direction is preserved but the magnitude is
reduced to $C$, i.e.
\[
\langle \nabla f(x_t),g_t\rangle = C\|\nabla f(x_t)\| .
\]
Thus the descent per iteration scales with $\min\{C,\|\nabla f(x_t)\|\}$,
which appears on the left‑hand side of Theorem~\ref{thm:main}.
\item \textbf{Comparison to vanilla SGD.}
If $C\ge \sup_t\|g_t^{\text{raw}}\|$ the clipping never activates and the bound
reduces to the classical $O(1/\sqrt{T})$ result for non‑convex SGD.
When gradients are heavy‑tailed, clipping yields the same $O(T^{-1/2})$ rate
while guaranteeing bounded updates.
\item \textbf{Hyper‑parameter Sensitivity.}
The bound depends linearly on $C$ through the term
$\min\{C,\|\nabla f\|\}$; a larger $C$ reduces the bias but may increase
the variance term $2L\eta\sigma^2$ if the raw gradients become unbounded.
\end{itemize}

%%%%%%%%%%%%%%%%%%%%%%%%%%%%%%%%%%%%%%%%%%%%%%%%%%%%%%%%%%%%%%%%%%%%%%%%%%%%
\section*{Preliminaries (Definitions and Lemmas)}
%%%%%%%%%%%%%%%%%%%%%%%%%%%%%%%%%%%%%%%%%%%%%%%%%%%%%%%%%%%%%%%%%%%%%%%%%%%%
\begin{lemma}[Descent Lemma for $L$‑smooth functions]\label{lem:descent}
For any $x$ and any vector $u$,
\[
f(x-u)\le f(x)-\langle\nabla f(x),u\rangle+\frac{L}{2}\|u\|^2 .
\]
\end{lemma}
\begin{proof}
Directly from the $L$‑smoothness inequality in Assumption~\ref{ass:smooth}.
\end{proof}

\begin{lemma}[Properties of the clipping operator]\label{lem:clip}
For any $v\in\mathbb{R}^d$ and threshold $C>0$ define $c(v)$ as in
\eqref{eq:clipped}. Then
\[
\|c(v)\|\le C,\qquad
\langle v, c(v)\rangle = \min\{C,\|v\|\}\,\|v\| .
\]
\end{lemma}
\begin{proof}
If $\|v\|\le C$, $c(v)=v$ and both statements hold trivially.
If $\|v\|>C$, $c(v)=C\,v/\|v\|$.
Hence $\|c(v)\|=C$ and
$\langle v,c(v)\rangle = C\|v\| = \min\{C,\|v\|\}\|v\|$.
\end{proof}

%%%%%%%%%%%%%%%%%%%%%%%%%%%%%%%%%%%%%%%%%%%%%%%%%%%%%%%%%%%%%%%%%%%%%%%%%%%%
\section*{Main Proof (Step‑by‑Step Derivation)}
%%%%%%%%%%%%%%%%%%%%%%%%%%%%%%%%%%%%%%%%%%%%%%%%%%%%%%%%%%%%%%%%%%%%%%%%%%%%
We prove Theorem~\ref{thm:main} by analysing one iteration and then
telescoping the inequality.

\noindent\textbf{Step 1 (Apply Descent Lemma).}
Using Lemma~\ref{lem:descent} with $u=\eta_t g_t$,
\[
f(x_{t+1}) \le f(x_t)-\eta_t\langle\nabla f(x_t),g_t\rangle
+\frac{L\eta_t^2}{2}\|g_t\|^2 . \tag{1}
\]

\noindent\textbf{Step 2 (Insert clipping identities).}
From Lemma~\ref{lem:clip} we have
\[
\langle\nabla f(x_t),g_t\rangle
=\min\{C,\|\nabla f(x_t)\|\}\,\|\nabla f(x_t)\| ,\qquad
\|g_t\|\le C . \tag{2}
\]

\noindent\textbf{Step 3 (Take conditional expectation).}
Condition on $x_t$ and use unbiasedness (Assumption~\ref{ass:unbiased}) and
variance bound (Assumption~\ref{ass:variance}):
\[
\mathbb{E}\!\bigl[g_t\mid x_t\bigr]=\nabla f(x_t),\qquad
\mathbb{E}\!\bigl[\|g_t\|^2\mid x_t\bigr]
= \|\nabla f(x_t)\|^2+\mathbb{E}\!\bigl[\|g_t-g_t^{\text{raw}}\|^2\mid x_t\bigr]
\le \|\nabla f(x_t)\|^2+\sigma^2 . \tag{3}
\]
(Note that clipping is a deterministic function of $g_t^{\text{raw}}$, thus
the variance bound still holds.)

\noindent\textbf{Step 4 (Combine (1)–(3)).}
Taking expectation of (1) conditioned on $x_t$ and applying (2)–(3) gives
\[
\begin{aligned}
\mathbb{E}\!\bigl[f(x_{t+1})\mid x_t\bigr]
&\le f(x_t)-\eta_t\,
\underbrace{\mathbb{E}\!\bigl[\langle\nabla f(x_t),g_t\rangle\mid x_t\bigr]}_{\text{use (2)}}\\
&\qquad
+\frac{L\eta_t^2}{2}\,
\underbrace{\mathbb{E}\!\bigl[\|g_t\|^2\mid x_t\bigr]}_{\le \|\nabla f(x_t)\|^2+\sigma^2}
\\[4pt]
&\le f(x_t)-\eta_t\,
\min\{C,\|\nabla f(x_t)\|\}\,\|\nabla f(x_t)\|
+\frac{L\eta_t^2}{2}\bigl(\|\nabla f(x_t)\|^2+\sigma^2\bigr). \tag{4}
\end{aligned}
\]

\noindent\textbf{Step 5 (Rearrange to isolate the min‑term).}
Bring the $\|\nabla f(x_t)\|^2$ term to the left‑hand side:
\[
\begin{aligned}
\eta_t\,
\min\{C,\|\nabla f(x_t)\|\}\,\|\nabla f(x_t)\|
&\le f(x_t)-\mathbb{E}[f(x_{t+1})\mid x_t]
+\frac{L\eta_t^2}{2}\|\nabla f(x_t)\|^2
+\frac{L\eta_t^2}{2}\sigma^2 .
\end{aligned}
\tag{5}
\]

\noindent\textbf{Step 6 (Use $\eta_t\le 1/L$).}
Because $\eta_t\le 1/L$, we have $L\eta_t^2/2\le \eta_t/2$.
Thus
\[
\frac{L\eta_t^2}{2}\|\nabla f(x_t)\|^2\le \frac{\eta_t}{2}\|\nabla f(x_t)\|^2 .
\tag{6}
\]

\noindent\textbf{Step 7 (Bound the product by the minimum).}
For any $a\ge0$ and $C>0$,
\[
\min\{C,a\}\,a \;\ge\; \frac12\min\{C^2,a^2\}.
\]
Indeed,
\begin{itemize}
\item If $a\le C$, then LHS$=a^2\ge \tfrac12 a^2$;
\item If $a>C$, then LHS$=Ca\ge \tfrac12 C^2$.
\end{itemize}
Hence
\[
\eta_t\,\min\{C,\|\nabla f(x_t)\|\}\,\|\nabla f(x_t)\|
\;\ge\;
\frac{\eta_t}{2}\,\min\{C^2,\|\nabla f(x_t)\|^2\}. \tag{7}
\]

\noindent\textbf{Step 8 (Plug (6) and (7) into (5).}
Using (6) to replace the positive $\tfrac{L\eta_t^2}{2}\|\nabla f\|^2$
by $\tfrac{\eta_t}{2}\|\nabla f\|^2$ and then (7) we obtain
\[
\frac{\eta_t}{2}\,\min\{C^2,\|\nabla f(x_t)\|^2\}
\;\le\;
f(x_t)-\mathbb{E}[f(x_{t+1})\mid x_t]
+\frac{\eta_t}{2}\|\nabla f(x_t)\|^2
+\frac{L\eta_t^2}{2}\sigma^2 .
\tag{8}
\]

\noindent\textbf{Step 9 (Cancel the $\tfrac{\eta_t}{2}\|\nabla f\|^2$ terms).}
Subtract $\frac{\eta_t}{2}\|\nabla f(x_t)\|^2$ from both sides:
\[
\frac{\eta_t}{2}\,\bigl(\min\{C^2,\|\nabla f(x_t)\|^2\}
-\|\nabla f(x_t)\|^2\bigr)
\le f(x_t)-\mathbb{E}[f(x_{t+1})\mid x_t]
+\frac{L\eta_t^2}{2}\sigma^2 .
\tag{9}
\]
Observe that $\min\{C^2,\|\nabla f\|^2\}-\|\nabla f\|^2\le 0$, therefore we can drop the
left‑hand side (it is non‑positive) and keep the inequality in the convenient
form
\[
\frac{\eta_t}{2}\,\min\{C^2,\|\nabla f(x_t)\|^2\}
\le f(x_t)-\mathbb{E}[f(x_{t+1})\mid x_t]
+\frac{L\eta_t^2}{2}\sigma^2 .
\tag{10}
\]

\noindent\textbf{Step 10 (Take full expectation).}
Let $\mathbb{E}_t[\cdot]=\mathbb{E}[\cdot\mid x_1]$ denote expectation over the
first $t$ random draws. Taking expectation of (10) yields
\[
\frac{\eta_t}{2}\,
\mathbb{E}\!\bigl[\min\{C^2,\|\nabla f(x_t)\|^2\}\bigr]
\le
\mathbb{E}\!\bigl[f(x_t)\bigr]-\mathbb{E}\!\bigl[f(x_{t+1})\bigr]
+\frac{L\eta_t^2}{2}\sigma^2 .
\tag{11}
\]

\noindent\textbf{Step 11 (Sum over $t=1,\dots,T$).}
Summing (11) telescopically,
\[
\frac{1}{2}\sum_{t=1}^{T}\eta_t
\mathbb{E}\!\bigl[\min\{C^2,\|\nabla f(x_t)\|^2\}\bigr]
\le
f(x_1)-\mathbb{E}\!\bigl[f(x_{T+1})\bigr]
+\frac{L\sigma^2}{2}\sum_{t=1}^{T}\eta_t^2 .
\tag{12}
\]

\noindent\textbf{Step 12 (Insert constant stepsize).}
With $\eta_t\equiv\eta\le 1/L$,
\[
\sum_{t=1}^{T}\eta_t = \eta T,\qquad
\sum_{t=1}^{T}\eta_t^2 = \eta^2 T .
\]
Using $f(x_{T+1})\ge f_{\inf}$ (Assumption~\ref{ass:lb}) we obtain
\[
\frac{\eta T}{2}\,
\frac{1}{T}\sum_{t=1}^{T}
\mathbb{E}\!\bigl[\min\{C^2,\|\nabla f(x_t)\|^2\}\bigr]
\le
f(x_1)-f_{\inf}
+\frac{L\eta^2\sigma^2}{2} T .
\tag{13}
\]

\noindent\textbf{Step 13 (Divide by $\eta T/2$).}
Rearranging (13) gives
\[
\frac{1}{T}\sum_{t=1}^{T}
\mathbb{E}\!\bigl[\min\{C^2,\|\nabla f(x_t)\|^2\}\bigr]
\le
\frac{2\bigl(f(x_1)-f_{\inf}\bigr)}{\eta T}
+2L\eta\,\sigma^2 .
\tag{14}
\]

\noindent\textbf{Step 14 (Replace $C^2$ by $C$).}
Since $\min\{C^2,a^2\}= \bigl(\min\{C,a\}\bigr)^2$, (14) is precisely the bound
stated in Theorem~\ref{thm:main}. \hfill $\square$

\section*{Rate Derivation (With Constants)}
From (14) we have
\[
\underbrace{\frac{1}{T}\sum_{t=1}^{T}
\mathbb{E}\!\bigl[\min\{C,\|\nabla f(x_t)\|\}^2\bigr]}_{\text{average clipped gradient squared}}
\le
\frac{2\Delta}{\eta T}+2L\eta\sigma^2,
\qquad
\Delta:=f(x_1)-f_{\inf}.
\]
Choosing $\eta = \sqrt{\frac{2\Delta}{L\sigma^2}}\;T^{-1/2}$ (which respects
$\eta\le 1/L$ for sufficiently large $T$) balances the two terms and yields
\[
\frac{1}{T}\sum_{t=1}^{T}
\mathbb{E}\!\bigl[\min\{C,\|\nabla f(x_t)\|\}^2\bigr]
\le
4\sqrt{2L\Delta}\,\sigma\, T^{-1/2}.
\]
Thus the algorithm attains the standard $O(T^{-1/2})$ rate for finding an
approximately stationary point, measured in terms of the clipped gradient.

\section*{Special Cases}
\begin{enumerate}[leftmargin=*]
\item \textbf{No clipping ($C\ge \sup_t\|g_t^{\text{raw}}\|$).}
Then $g_t=g_t^{\text{raw}}$, $\min\{C,\|\nabla f\|\}=\|\nabla f\|$, and (14) reduces
to the classic non‑convex SGD bound
\[
\frac{1}{T}\sum_{t=1}^{T}\mathbb{E}\bigl[\|\nabla f(x_t)\|^2\bigr]
\le
\frac{2\Delta}{\eta T}+2L\eta\sigma^2 .
\]

\item \textbf{Heavy‑tailed gradients.}
If the raw gradient distribution has infinite variance but the clipping
threshold $C$ is finite, then $\|g_t\|\le C$ guarantees that all moments of
$g_t$ exist, and the variance term $\sigma^2$ in the bound can be replaced by
the (finite) variance of the clipped gradients, yielding the same $O(T^{-1/2})$
rate while avoiding divergence.

\item \textbf{Deterministic setting ($\sigma=0$).}
When exact gradients are available, the bound simplifies to
\[
\frac{1}{T}\sum_{t=1}^{T}
\mathbb{E}\!\bigl[\min\{C,\|\nabla f(x_t)\|\}^2\bigr]
\le
\frac{2\Delta}{\eta T},
\]
and choosing a constant stepsize $\eta=1/L$ gives a $O(1/T)$ decay of the
average clipped gradient norm.
\end{enumerate}

\section*{Discussion}
Clipped SGD preserves the direction of the true gradient while capping its
magnitude. The convergence analysis above shows that the clipping bias is
exactly captured by the $\min\{C,\|\nabla f\|\}$ term; consequently the algorithm
inherits the optimal $O(T^{-1/2})$ rate for non‑convex stochastic optimisation.
Moreover, the boundedness of the updates makes the method robust to outlier
gradients, a property that is not guaranteed for vanilla SGD.

\end{document}